
\section{Conclusion}
\label{conclude}

This paper introduced SECDA-DSE, a framework that extends the SECDA ecosystem with automated design space exploration for FPGA-based AI accelerators harnessing LLMs. 
SECDA-DSE combines two major components, the DSE Explorer and an intelligent LLM Stack, to guide the exploration of hardware accelerator architectures through reasoning and context driven decision making. 
SECDA-DSE enables the generation and refinement of hardware accelerator configurations based on workload requirements, device constraints, and architectural directives through retrieval-augmented generation. Furthermore, chain-of-thought reasoning, and parameter-efficient fine-tuning through LoRA enables efficient model adaptation by freezing the base model parameters and learning small low-rank matrices that specialize model to a new task, reducing computational and memory costs during fine-tuning. 
SECDA-DSE leverages LLMs to analyze hardware architecture trade-offs and iteratively improve generated accelerator designs through reinforcement guided by human-in-the-loop feedback. 
Our preliminary evaluation demonstrates that SECDA-DSE can successfully translate high-level natural language specifications into SECDA-native accelerator designs that meet HLS synthesis constraints, validating the feasibility of LLM-driven accelerator generation.

% SECDA-DSE focuses on harnessing the potential of LLM-driven reasoning with traditional hardware design workflows to reduce the complexity of accelerator development. 
Future work will focus on fully automating the design space exploration loop by integrating the DSE Explorer and LLM Stack through an MCP interface, expanding the evaluation, and exploring the effectiveness of multiple fine-tuned open-source LLMs across larger hardware design datasets.